\subsection*{Висновки до другого розділу}
\addcontentsline{toc}{subsection}{Висновки до другого розділу}

Комплексний аналіз стану формування ІКК в українській дизайн-освіті на двох рівнях (загальнонаціональному та інституційному) з міжнародним бенчмаркінгом дозволяє сформулювати такі висновки.

1. Загальнонаціональний обчислювальний аналіз 122 акредитованих програм спеціальності 022~Дизайн (3\,043 освітні компоненти, 1\,515 унікальних назв дисциплін) виявив системну прогалину у формуванні ІКК. При задовільному покритті традиційних цифрових інструментів (71,4\%), компоненти <<Основи ШІ>>, <<Генеративний ШІ>> та <<Комп'ютерний зір>> мають \textbf{0\% покриття}. Загальний рівень інтеграції ШІ-тематики становить лише 35,7\%. Мережевий аналіз (1\,240 вузлів, 7\,524 ребер, 10 спільнот Лувена) підтвердив, що цифрові/ШІ-дисципліни не утворюють самостійного кластера. Фактична відсутність вибіркових компонентів ($<$0,1\%) свідчить про ригідність навчальних планів.

2. Кейс-стаді трьох програм КДПУ (ГД, ДО, ДС) із тріангуляцією 6 джерел даних (8~версій ОПП, 132~силабуси, опитування 32~випускників, 9~роботодавців, 61~студента, висновки 3~експертних груп НАЗЯВО) виявив:
\begin{itemize}
    \item нисхідний тренд ІКК-ключових слів в ОПП (з~122 до 75--82 між 2018 та 2021~рр., $-$35\%);
    \item нульове покриття ШІ-тематики у 132 проаналізованих документах;
    \item незадоволеність випускників ІКТ-підготовкою (33,3\% у ГД, 15,4\% у ДО);
    \item пряму рекомендацію роботодавців щодо інтеграції ШІ;
    \item статистично значущі міжпрограмні відмінності у задоволеності матеріально-технічним ($H = 8{,}708$, $p = 0{,}013$) та інформаційним ($H = 6{,}019$, $p = 0{,}049$) забезпеченням.
\end{itemize}

3. Міжнародний бенчмаркінг 7 провідних програм дизайну (Великобританія, США, Фінляндія, Нідерланди, Австралія, Італія, Південна Корея) підтвердив суттєвий розрив: міжнародні програми мають від 1 до 4 ШІ-дисциплін, інтегрують промпт-інженерію, ШІ-етику та критичний дизайн, тоді як українські програми мають 0 таких дисциплін. Визначено 5 найкращих практик для адаптації: прогресивна інтеграція ШІ, збереження ручних навичок, структурована промпт-інженерія, ШІ-етика, портфоліо-оцінка.

Виявлені дані обґрунтовують необхідність розробки моделі трансформації навчальних програм для формування ІКК~--- предмет розділу~3.
